% ECONOMICS_PROBLEM: {"id": "C78-134", "title": "Does double primality reduce the maximum number of stable matchings?", "jel_code": "C78", "source_id": "OP-0218"}
% !TeX program = xelatex
\documentclass[11pt,a4paper]{article}
\usepackage[margin=23mm]{geometry}
\usepackage{fontspec}
\setmainfont{Latin Modern Roman}
\usepackage{amsmath,amssymb,mathtools}
\usepackage{xcolor,enumitem,booktabs,longtable,array,xurl,hyperref}
\definecolor{muted}{HTML}{53636C}
\hypersetup{colorlinks=true,urlcolor=blue,linkcolor=blue}
\setlength{\parindent}{0pt}
\setlength{\parskip}{5pt}
\newcommand{\R}{\mathbb R}
\newcommand{\E}{\mathbb E}
\newcommand{\N}{\mathbb N}
\newcommand{\ind}{\mathbf 1}
\newcommand{\Var}{\operatorname{Var}}
\newcommand{\Cov}{\operatorname{Cov}}
\newcommand{\supp}{\operatorname{supp}}
\DeclareMathOperator*{\argmin}{arg\,min}
\DeclareMathOperator*{\argmax}{arg\,max}
\newcommand{\entrymeta}[3]{{\sffamily\small\color{muted}\textbf{\texttt{#1}}\quad|\quad #2\par #3\par}}
\newcommand{\Problemref}[1]{\texttt{#1}}
\newcommand{\JELref}[2]{\texttt{#2} (JEL \texttt{#1})}
\begin{document}
\section*{Does double primality reduce the maximum number of stable matchings?}
\textbf{Problem ID:} \texttt{C78-134}\quad \textbf{JEL:} \texttt{C78}\par
% BEGIN ECONOMICS STATEMENT
\label{problem:OP-0218}
% GAME_THEORY_PROBLEM: {"local_id": "GT-088", "original_number": 88, "kind": "P", "entry_kind": "statement_only", "published": false, "review_required": true, "category": "game_theory"}
\label{game:GT-088}
{\sffamily\small\color{muted}
\noindent\textbf{ID:} GT-088. \textbf{Category:} Stable matching and extremal combinatorics.
\textbf{Status:} P; question E in the inspected source, with computational qualifications below.
\par}
\subsection*{Definitions and mathematical statement}
For an integer $n\geq 1$, an instance $I$ consists of disjoint labelled sets $P,R$, each of cardinality $n$, and a strict total preference order over the opposite set for each agent. A matching is a bijection $\mu:P\to R$. It is stable if there is no pair $(p,r)$ with
\[
 r\succ_p\mu(p)\quad\text{and}\quad p\succ_r\mu^{-1}(r).
\]
Write $\operatorname{Stab}(I)$ for its stable matchings.

An \emph{input-independent partition} consists of partitions
$P=\bigsqcup_{j=1}^tP_j$, $R=\bigsqcup_{j=1}^tR_j$ with $|P_j|=|R_j|>0$, such that every member of $P_j$ ranks all of $R_j$ ahead of $R\setminus R_j$, and every member of $R_j$ ranks all of $P_j$ ahead of $P\setminus P_j$. The instance is \emph{input-prime} if no such partition with $t\geq2$ exists.

The proposer-side deferred-acceptance process has states $(h,q)$: $h(r)\in P\cup\{\varnothing\}$ is the proposer held by $r$, and $q(p)$ is the number of proposals already made by $p$. Initially all holders are empty and $q=0$. An unheld proposer with an untried receiver proposes to the next receiver on its list; that receiver keeps its preferred proposer. Consider all states reachable under all legal serial schedules. Put
\[
 \mathcal F(I)=\left\{\{(p,j):0\leq j<q(p)\}: (h,q)\text{ is reachable}\right\},
 \qquad E(I)=\bigcup_{F\in\mathcal F(I)}F.
\]
The event system is \emph{behaviourally decomposable} if a partition
$E(I)=\bigsqcup_{j=1}^tE_j$, $t\geq2$, into nonempty parts satisfies
\[
 \mathcal F(I)=\left\{\bigsqcup_{j=1}^tF_j:
 F_j\in\{F\cap E_j:F\in\mathcal F(I)\}\right\}.
\]
Otherwise it is behaviourally prime. An instance is \emph{doubly prime} if it is input-prime and behaviourally prime. Define
\[
 S_n=\max_I|\operatorname{Stab}(I)|,\qquad
 S_n^{\mathrm{DP}}=\max_{I\text{ doubly prime}}|\operatorname{Stab}(I)|,
\]
with the second maximum taken as $0$ if its indexing class is empty. The question is whether
\[
 \forall n\geq1\ (n\neq4\Longrightarrow S_n^{\mathrm{DP}}=S_n).
\]

\subsection*{Short English statement}
Except at order four, can an instance with the largest possible stable set always be chosen to have neither kind of nontrivial decomposition?

\subsection*{Variants and scope boundaries}
The behavioural definition preserves labelled proposal events; a factorization of an uncoloured execution graph is a different condition. The definition depends on the designated proposing side. Ties, incomplete lists, and unequal sides are excluded. A common-receiver connectivity test is not substituted for the event-family definition.

\subsection*{Source relationship and status}
The source poses this extremal question in Section 7, question E. It reports $S_4^{\mathrm{DP}}=6$ versus $S_4=10$, equality at order three, and corresponding evidence at larger orders. In particular its order-six figure is a search maximum, not a proof of the unrestricted extremum. These source-reported computations do not prove the quantified assertion.

\subsection*{References}
\begingroup\footnotesize\raggedright
\noindent [S1] Yoshiteru Ishida, \emph{Prime factorisation of stable-matching instances: uniqueness, simultaneous products, and an exact census} (2026), arXiv:2610.05617v1. Inspected locators: Section 2 (reachable events and the two decompositions), Section 4 (double primality), Section 7, question E.
\url{https://arxiv.org/html/2610.05617v1}
\par\endgroup

\subsection*{Common conventions: Source mathematical conventions}

\textbf{Finite games.} For a finite set $A$, $\Delta(A)$ is its probability
simplex. Players choose independent mixed strategies in Nash equilibrium;
correlation is allowed only when explicitly part of the solution concept.
In a game with payoff functions $u_i$ and strategy profile $x$, an additive
$\varepsilon$ Nash equilibrium satisfies
\[
 u_i(x)\geq u_i(x_i',x_{-i})-\varepsilon
 \quad\text{for every player $i$ and every admissible deviation $x_i'$}.
\]
For numerical approximation, payoffs are normalized as locally specified.
An $\varepsilon$ payoff residual is distinct from distance at most
$\varepsilon$ to the set of exact equilibria. Point convergence, convergence
of empirical play, and convergence in distance to an equilibrium set are
also distinct.

\textbf{Computational representations.} Unless stated otherwise, a complexity
question takes rational data with binary encoding; polynomial time is measured
in total bit length. A fixed-accuracy polynomial algorithm is not necessarily
polynomial in $1/\varepsilon$, and neither is necessarily polynomial in
$\log(1/\varepsilon)$. Succinct games and value, demand, or sampling oracles
must declare their interfaces. Exact real solutions require an explicit
representation; an unspecified list of arbitrary real numbers is not a
finite computational output. A claimed algorithmic barrier is meaningful
only for its specified model and reductions.

\textbf{Dynamic games.} Histories, information, observation, and allowable
randomization are part of the model. A stationary strategy depends only on
the current state; a positional strategy in a deterministic graph game
chooses one action at each controlled vertex. Nash equilibrium at an initial
state and equilibrium after every history are different requirements.
An expectation of a pathwise $\liminf$ payoff, a $\liminf$ of expected averages,
a limiting expected average, and a uniform finite-horizon equilibrium are
different criteria. None is silently substituted for another.

\textbf{Allocation and mechanisms.} A complete allocation assigns every item
exactly once unless otherwise stated. Zero-valued goods, negative values,
capacity constraints, feasibility, lotteries, and copies of goods affect
fairness definitions and are specified locally. Dominant-strategy incentive
compatibility, Bayesian incentive compatibility, universal truthfulness,
and truthfulness in expectation impose different inequalities. Whenever
an objective ratio has zero denominator, its convention is stated or those
instances are excluded.

\textbf{Infinite-dimensional models.} Expectations, integrals, stochastic
equations, and optimization objectives require their local regularity,
integrability, filtrations, and admissible domains. A backward--forward
mean-field system is a boundary-value problem; it is not an initial-value
semiflow without an additional construction. Long-time, large-population,
and vanishing-noise limits require an order of limits and a convergence
topology. If a minimum need not be attained, the objective is its infimum
or supremum and the approximation requirement is stated separately.
% END ECONOMICS STATEMENT
\end{document}
